\section{Exact-Action Risk: A Type-Constrained Spectral Characterization}
\label{sec:boundaries}
The optimum $\eta$ answers how accurately unrelated linear maps can predict the transformed observations. It does not tell us how much accuracy must be sacrificed to use one exact orthogonal action. We show what the scalar regression risk and the untyped spectrum of the orbit Gram matrix leave unresolved, then construct a spectral quantity that bounds $L^\star$. Throughout this section the coordinates are fixed, states and group elements are averaged uniformly, and the uncentered second moment is positive definite, as specified in \cref{sec:setup}.

\subsection{Equal regression risks can hide different exact-action risks}

We use the common loss in \cref{eq:class-loss-main}, with a
single template and observation dimension $r=1$.
Let $G=S=C_2\times C_2\times C_7$, with states
$s=(a,b,k)$ and group elements $g=(u,v,w)$.
The prescribed action is componentwise addition:
\[
 g\cdot(a,b,k)=(a+u,b+v,k+w),
\]
where the first two coordinates are taken modulo $2$ and
the third modulo $7$.
We average uniformly over all $28$ states and all $28$
group elements, including the identity.
On this same state action, consider two scalar observation
systems:
\[
 x_A(a,b,k)=\frac{(-1)^a+(-1)^b}{\sqrt{2}},
 \qquad
 x_B(a,b,k)=\sqrt{2}\cos(2\pi k/7).
\]
Both have zero mean and unit second moment over the complete
state set, so they already satisfy the whitening condition.
Centering this single complete template leaves both
observations unchanged.

For $J\in\{A,B\}$, the optimal independent scalar regression
coefficient is the cross-moment
$m_J(g)=\mathbb E_s[x_J(s)x_J(g\cdot s)]$.
Here,
\[
 m_A(u,v,w)=\frac{(-1)^u+(-1)^v}{2},
 \qquad
 m_B(u,v,w)=\cos(2\pi w/7).
\]
Both satisfy $\mathbb E_g[m_J(g)^2]=1/2$, giving
\[
 \eta_A=\eta_B=\frac12.
\]
Their orbit Gram matrices, which record cross-moments among
all group-transformed observations with the normalization
in \cref{eq:orbit-gram-main}, also have the same nonzero
eigenvalues:
\[
 \operatorname{spec}_{\ne 0}(K_A)
 =
 \operatorname{spec}_{\ne 0}(K_B)
 =
 \left\{\frac12,\frac12\right\}.
\]

The distinction appears in the invariant subspaces spanned
by the transformed observation functions.
As $g$ varies, the functions $s\mapsto x_J(g\cdot s)$ span
\[
 \begin{aligned}
 V_A
 &=\operatorname{span}_{\mathbb R}
   \left\{(-1)^a,\,(-1)^b\right\},\\
 V_B
 &=\operatorname{span}_{\mathbb R}
   \left\{\cos(2\pi k/7),\,\sin(2\pi k/7)\right\}.
 \end{aligned}
\]
These are two-dimensional spaces of functions on $S$;
each observation itself remains scalar.
The space $V_A$ splits into two distinct real
one-dimensional characters.
The space $V_B$ is a real irreducible two-dimensional
rotation representation and has no invariant real line.
Thus a one-dimensional action can represent either
character in $V_A$, but cannot represent one coordinate
of the rotation plane in $V_B$ as an invariant component.

More explicitly, an exact orthogonal action in one real
dimension is a character $\chi:G\to\{+1,-1\}$.
Every such character is trivial on the $C_7$ factor,
because its generator has odd order.
For observation system $J$, its loss is
\[
 \mathcal L_J(\chi)
 =
 2-2\mathbb E_g[\chi(g)m_J(g)].
\]
The largest correlation is $1/2$ for $A$, attained by
$\chi(u,v,w)=(-1)^u$ or $\chi(u,v,w)=(-1)^v$.
For $B$, the correlation is zero for every admissible
character.
Consequently,
\[
 \boxed{L_A^\star=1,\qquad L_B^\star=2.}
\]

The independently orthogonal comparators also distinguish
the two systems:
$L_{\mathrm{orth},A}=L_A^\star=1$, whereas
$L_{\mathrm{orth},B}<L_B^\star=2$.
Enforcing the group law therefore adds no further cost
in $A$ and a strictly positive cost in $B$.
The complete calculation, including the separation of
orthogonality and group-law costs, is given in
\cref{sp:counterexample}.

This example shows that the independent linear optimum
and the ordinary orbit-Gram eigenvalues do not determine
the exact-action optimum.
The complete state-labeled observations distinguish the
two systems.
The information omitted by these summaries is how the
corresponding subspaces transform under the group and
which of them fit within the prescribed real dimension.

\subsection{Measuring energy under an invariant-dimension budget}

We now retain the group-action information omitted by the
ordinary spectrum.
Orbit Gram matrices provide a finite formulation of
invariant approximation \citep{barbieri2020approximation}.
We use the whitened representation rows $x_{i,t}$ from
the preceding setup, with
$\mathbb E_{t,i}[x_{i,t}^{\top}x_{i,t}]=I_r$.
Expectations use the fixed template distribution and an
independent uniform state.
For each template--state pair, concatenate the transformed
observations:
\begin{equation}
 v(t,i)
 =
 |G|^{-1/2}\bigl(x_{g(i),t}\bigr)_{g\in G}
 \in\R^{1\times r|G|},
 \qquad
 K=\mathbb E_{t,i}[v^\top v].
 \label{eq:orbit-gram-main}
\end{equation}
Each group element contributes one block, including the
identity and elements that fix the source state.
Repeated targets produced by different group elements
are retained.
For example, in the square-corner task, the eight elements
send each source corner to each of the four targets twice.

The matrix $K$ is positive semidefinite and satisfies
$\operatorname{tr}K=r$ under the unit-second-moment
condition.
Writing $x_g(i,t)=x_{g(i),t}$, its blocks are
\[
 K_{g,h}
 =
 |G|^{-1}\mathbb E_{t,i}[x_g^\top x_h].
\]
In particular,
$|G|K_{e,g}=\mathbb E_{t,i}[x^\top x_g]$
is the cross-moment used in the common prediction loss.
This construction uses complete labeled observation
blocks.
The \emph{ordinary orbit-Gram spectrum} consists of the
eigenvalues of $K$, counted with multiplicity, without
information about how their eigenspaces transform under
the group.

The known group multiplication also defines permutations
of the block indices.
For a block column vector
$a=(a_k)_{k\in G}$, with $a_k\in\R^{r\times1}$, set
\[
 (\Lambda_g a)_k=a_{kg^{-1}}.
\]
These orthogonal block permutations satisfy
$\Lambda_g\Lambda_h=\Lambda_{hg}$, consistently with the
reversed multiplication order of the row-action convention.
Uniform state averaging gives
$\Lambda_g^\top K\Lambda_g=K$, and hence
$K\Lambda_g=\Lambda_gK$.
A subspace is invariant if every $\Lambda_g$ maps it
into itself.

We seek the invariant subspace that captures the most
orbit energy using at most $r$ real dimensions.
For projections
$Q\in\R^{r|G|\times r|G|}$, define
\begin{equation}
 \begin{aligned}
 C_r(K)
 &=
 \max_Q\operatorname{tr}(KQ),
 \qquad
 \gamma_r=\frac{C_r(K)}{r},\\
 \text{subject to}\quad
 &Q^2=Q=Q^\top,\qquad
   Q\Lambda_g=\Lambda_gQ\quad(\forall g\in G),\\
 &\operatorname{rank}Q\le r,\qquad
   \operatorname{range}Q\subseteq\operatorname{range}K.
 \end{aligned}
 \label{eq:budget-main}
\end{equation}
The projection acts in the concatenated orbit coordinate
space.
Commutation enforces invariance, and the support restriction
limits selection to directions carrying observation energy.
The rank inequality permits unused dimensions when no
additional complete invariant component fits.
Since $\operatorname{tr}K=r$, the normalized budget
$\gamma_r\in[0,1]$ is the fraction of total energy captured.
The dependence of $C_r(K)$ on the known block action
$\Lambda$ is implicit in this notation.

An irreducible real representation has no nonzero proper
invariant subspace; its type specifies its equivalence
class \citep{serre1977linear}.
Because $K$ commutes with the block action, its eigenspaces
can be decomposed into real irreducible copies.
A copy of type $\tau$ occupies $d_\tau$ real dimensions
and carries an eigenvalue $\lambda$ repeated $d_\tau$
times, contributing energy $d_\tau\lambda$.
The optimization selects complete copies whose dimensions
sum to at most $r$.
The corresponding integer optimization is given in
\cref{sp:integer}.
For the preceding example,
\[
 C_1(K_A)=\frac12,
 \qquad
 C_1(K_B)=0.
\]
System $A$ permits selection of one character.
System $B$ has only an irreducible rotation plane in its
support, so a one-dimensional budget admits only the
zero projection.

Ordinary spectral truncation can split such an irreducible
component and therefore overestimate the feasible energy.
Define
\[
 \Delta_{\mathrm{type}}
 :=
 \sum_{j=1}^{r}\lambda_j(K)-C_r(K)
 \ge0,
 \qquad
 \lambda_1(K)\ge\lambda_2(K)\ge\cdots,
\]
where eigenvalues are counted with multiplicity, including
zeros.
This quantity measures the budget overestimate from
ignoring the invariant-subspace constraint.
It is distinct from a prediction risk or a group-law cost.
Under the same whitening and averaging contract, the
optimal independent regression risk satisfies
\[
 \eta=1-\frac{\operatorname{tr}(K^2)}{r}.
\]
Thus this optimal risk depends only on the ordinary
eigenvalues, whereas $C_r(K)$ also uses their real
representation types.
The regression matrices themselves depend on the
individual cross-moments.
The finite formulation and its invariant-approximation
equivalent are established in
\cref{sp:budget,sp:observable}.

For a fixed abstract group, coordinate dimension, and
group-element indexing, $\Lambda$ is determined by the
group multiplication.
Changing the registration changes the target
correspondences and can change $K$ while leaving
$\Lambda$ fixed.
We compare candidate budgets on the same representation
bank and in the same metric.
Their ordering describes spectral compatibility with
the candidate correspondences; an ordering of the
exact-action optima requires separate justification
from the corresponding risk bounds.

The budget includes both trivial and nontrivial types.
For a transitive action on a complete, uniformly weighted
state set, the invariant component corresponds to the
within-template state mean
(\cref{sp:transitive-trivial-energy}).
Removing this component while retaining the original
whitening deletes the trivial block of $K$ and preserves
its nontrivial blocks.
It also releases dimensions previously allocated to
trivial copies, so the budget must be optimized again.
Consequently, a common trivial spectrum across candidates
does not imply that their rankings are unchanged by
centering.

The centering analysis in
\cref{app:template-centering-sensitivity}
distinguishes this fixed-metric operation from whitening
the centered observations by their own second moment.
The original and rewhitened centered treatments satisfy
the unit-second-moment condition after their respective
positive-definiteness checks.
For centered observations retaining the original whitening,
$C_r(K)/r$ instead measures captured energy relative to
the original total energy.
When $\operatorname{tr}K>0$, the fraction of remaining
energy captured is $C_r(K)/\operatorname{tr}K$.
This treatment generally lacks unit second moment, so
neither the regression identity above nor the subsequent
unit-second-moment risk bounds apply directly.
Its reported budgets and registration ranks assess
sensitivity to removal of the state-invariant component.

\subsection{From the budget to risk and constraint costs}

The spectral budget bounds the minimum prediction risk over
the exact orthogonal-action class.
All class optima and spectral quantities below use the same
observation distribution, prescribed state action, averaging
weights, and whitening metric.

\begin{spectraltheorem}[Risk of the entire exact orthogonal action class]
\label{sp:main}
Let a finite group $G$ act by permutations on a finite state
set $S$, with complete fixed-coordinate observations
$z_{i,t}\in\R^{1\times r}$, $r\ge1$.
Expectations average over $t\sim P_T$ and an independent
uniform state $i\in S$.
Group averages include every $g\in G$, with uniform weights,
including the identity and fixed-point pairs.
Assume that the uncentered second moment
$\Sigma=\mathbb E_{t,i}[z_{i,t}^{\top}z_{i,t}]$
is finite and positive definite, and set
$x_{i,t}=z_{i,t}\Sigma^{-1/2}$.
Let $L^\star$ be the minimum loss in
\cref{eq:class-loss-main} over $\mathcal R_r(G)$, where
$\rho_g\in O(r)$, $\rho_e=I_r$, and
$\rho_g\rho_h=\rho_{hg}$ for every $g,h\in G$.
Let $\gamma_r$ be the normalized invariant-dimension budget
in \cref{eq:budget-main}.
Then
\begin{equation}
 \boxed{2(1-\gamma_r)\le L^\star\le2(1-\gamma_r^2).}
 \label{eq:main-risk-bound}
\end{equation}
Every member of $\mathcal R_r(G)$ satisfies the lower bound,
and some member of this same $r$-dimensional class attains
a loss no greater than the upper bound.
\end{spectraltheorem}

The state action may be nonfree or nonfaithful, and
$\mathcal R_r(G)$ permits nonfaithful representations.
The upper endpoint guarantees the existence of a suitable
action on the specified observation distribution.
It does not guarantee the performance of an arbitrary
constructed action or its transfer to new templates.
When $\gamma_r<1$, the ratio of the two endpoints is
$1+\gamma_r\le2$; at $\gamma_r=1$, both endpoints are zero.
This factor-of-two statement concerns the bounds on
$L^\star$ itself and need not hold after subtracting a
comparison-class optimum.

The proof relates prediction to invariant approximation.
Averaging the observation map against a candidate exact
action projects it onto the space of intertwining maps,
which respect the group actions on their domain and range.
The image of such a map is an invariant subspace of
dimension at most $r$, and the spectral budget bounds its
alignment with the observations.
Conversely, a subspace attaining the budget, orthogonally
aligned with the observations, yields an admissible action
with controlled risk.
If its dimension is less than $r$, the construction uses
auxiliary directions carrying the trivial action in the
proof's ambient function space.
The observed coordinates in those directions are zero,
and the final predictor matrices remain $r\times r$.
The complete argument is given in \cref{sp:main-proof}.

\paragraph{The combined cost of exact laws and orthogonality.}
Subtracting the independent linear optimum from the risk
bounds gives
\begin{equation}
 p^-:=\max\{0,2(1-\gamma_r)-\eta\}
 \le L^\star-\eta\le
 p^+:=2(1-\gamma_r^2)-\eta.
 \label{eq:risk-premium-main}
\end{equation}
These endpoints bound the joint cost of imposing
orthogonality and exact group laws under the common loss.
The spectral and regression identities also imply
\begin{equation}
 \gamma_r\le\sqrt{1-\eta},
 \qquad
 p^-\ge(1-\sqrt{1-\eta})^2.
 \label{eq:positivity-main}
\end{equation}
Thus $\eta>0$ already entails a positive lower endpoint
in exact arithmetic.
Positivity alone therefore establishes neither a substantial
premium nor a positive additional group-law cost.

\paragraph{How much does orthogonality alone explain?}
Let $M_g=\mathbb E_{t,i}[x^\top x_g]$.
Orthogonal Procrustes optimization
\citep{schonemann1966procrustes} gives
\[
 L_{\mathrm{orth}}
 =
 2-\frac{2}{r|G|}\sum_{g\in G}\|M_g\|_*,
 \qquad
 \beta=L_{\mathrm{orth}}-\eta,
\]
where the nuclear norm sums the singular values.
Each group element chooses its own matrix in $O(r)$,
including possible reflections, without a shared group-law
constraint.
If $\sigma_{g,1},\ldots,\sigma_{g,r}$ are all the singular
values of $M_g$, including zeros, then
\[
 \beta
 =
 \frac{1}{r|G|}
 \sum_{g\in G}\sum_{j=1}^{r}(1-\sigma_{g,j})^2.
\]
Under the whitening assumptions, these singular values lie
in $[0,1]$.
In particular, $\eta>0$ already implies a positive
orthogonality cost.

The joint cost separates as
\[
 L^\star-\eta
 =
 \beta+(L^\star-L_{\mathrm{orth}}).
\]
The additional group-law cost satisfies
\[
 \begin{aligned}
 \ell^-
 &:=
 \max\{0,2(1-\gamma_r)-L_{\mathrm{orth}}\},\\
 \ell^+
 &:=
 2(1-\gamma_r^2)-L_{\mathrm{orth}},\\
 \ell^-
 &\le L^\star-L_{\mathrm{orth}}\le\ell^+.
 \end{aligned}
\]
The cost $\beta$ is directly computable from the two
comparison-class optima.
The remaining cost is unknown and is bounded by
$[\ell^-,\ell^+]$.
In exact arithmetic, $\ell^->0$ establishes a positive
additional cost, whereas $\ell^-=0$ by itself leaves that question
unresolved.
Although the true costs add, their lower endpoints satisfy
\[
 \beta+\ell^-=\max\{\beta,p^-\},
\]
which can exceed $p^-$.
A ratio such as $\beta/p^-$ therefore does not measure
the fraction of the true joint cost due to orthogonality.
All these comparisons use class optima under the same loss.
Reported errors of training-fitted OLS, Ridge, Joint, or
transferred witnesses cannot replace them without establishing
both the same evaluation loss and attainment of the relevant
class optimum.
The derivations are given in
\cref{sp:premium-positivity,sp:orthogonal-comparator}.

The additional group-law cost is zero exactly when the
independently optimal orthogonal maps admit a common
choice satisfying the row-action laws.
If every $M_g$ is nonsingular, its orthogonal polar factor
is unique, so compatibility reduces to checking that family.
With singular cross-moments, an arbitrary SVD completion
may fail the laws even when another optimal choice satisfies
them (\cref{sp:orthogonal-compatibility}).
This criterion concerns exact compatibility; a small
floating-point law residual alone certifies neither exact
compatibility nor a strictly positive additional cost.

The two observation systems in the preceding example
illustrate this separation.
With
$a_7=\frac17\sum_{k=0}^{6}|\cos(2\pi k/7)|$, their costs are
\[
 \begin{array}{c|ccc}
  & L_{\mathrm{orth}}
  & L_{\mathrm{orth}}-\eta
  & L^\star-L_{\mathrm{orth}}\\
  \hline
  A & 1 & 1/2 & 0\\
  B & 2-2a_7 & 3/2-2a_7 & 2a_7
 \end{array}
\]
The systems share $\eta=1/2$ and the ordinary orbit-Gram
spectrum, but have different independently orthogonal optima.
System $A$ admits a compatible optimal choice despite
singular cross-moments.
In $B$, the unique polar signs violate the cycle law,
giving a positive additional cost.
These are exact calculations for the analytical example;
\cref{sp:counterexample} gives the derivation.

\paragraph{An identity reference in the same metric.}
The exact action assigning $\rho_g=I_r$ to every group
element predicts no change under any transformation.
Under the unit-second-moment condition, its risk is
\begin{equation}
 \begin{aligned}
 L(I)
 &=
 \mathcal L((I_r)_{g\in G})
 =
 2(1-e_{\mathrm{triv}}),\\
 e_{\mathrm{triv}}
 &=
 \frac1r\mathbb E_{t,i}
 \left\|
   \frac1{|G|}\sum_{g\in G}x_g
 \right\|_2^2.
 \end{aligned}
 \label{eq:identity-reference-main}
\end{equation}
This evaluates the identity family over all group elements.
The group average extracts the state-invariant component,
which is the within-template state mean for the transitive
weekday and corner actions.
The quantity $e_{\mathrm{triv}}$ counts all trivial-type
energy, normalized by $r$.
The selected share $b_{\mathrm{triv}}$ counts only the part
retained by a maximizing dimension budget.
On the same bank and in the same metric, all transitive
candidate images have the same $e_{\mathrm{triv}}$ and
$L(I)$, while the selected share can vary with the dimension
allocation.
Since the identity family is admissible,
$L^\star\le L(I)$
(\cref{sp:identity-reference,sp:transitive-trivial-energy}).

The normalization matters for the centering analysis in
\cref{app:template-centering-sensitivity}.
For its transitive actions, within-template centering
removes the invariant component.
When the centered second moment is positive definite,
rewhitening gives $e_{\mathrm{triv}}=0$ and $L(I)=2$.
Expanding the squared loss for centered observations retaining
the original whitening gives an identity-family risk of
$2E^{(\mathrm{fixed})}/r$ under the $1/r$ normalization, where
$E^{(\mathrm{fixed})}
=\mathbb E_{t,i}\|x^{(\mathrm{fixed})}_{i,t}\|_2^2$.
That residual energy generally differs from $r$, so the
unit-second-moment formula cannot be applied directly to
this treatment.

\paragraph{What the calculation establishes.}
For a complete finite bank, the theorem gives deterministic
mathematical bounds on that bank's exact orthogonal-action
optimum.
The reported numerical endpoints approximate these bounds
under the stated rank and whitening checks.
A budget near one guarantees that some exact action has
low risk in the declared metric, and $L^\star=0$ holds
exactly when $\gamma_r=1$.

Zero risk need not select a unique registration.
For a fixed set of $n$ affinely independent state vectors,
retain their full $(n-1)$-dimensional centered span and
whiten using the uniformly weighted state second moment.
Every state permutation then has an exact orthogonal
realization (\cref{prop:finite}).
In such a shared-code construction, multiple registrations
can have zero risk.
Separate realizability within individual templates,
however, does not ensure a common action across templates.
For fitted linear maps, composition residuals also depend
on prediction residuals and operator norms
(\cref{prop:propagation}).

The class studied here consists of exact orthogonal actions
in the specified whitened coordinates.
Ordinary Ridge and soft-penalty Joint fits need not belong
to this class.
Two whitenings of the same second moment differ by an
orthogonal coordinate change and preserve orthogonality.
Changing to a metric defined by a different bank or
distribution preserves group laws under conjugation but
can change Euclidean orthogonality.
The transferred witness uses frozen training-derived
matrices and training whitening on held-out templates.
Its test loss is a prediction score in that training metric,
whereas the finite-bank bounds use the evaluated bank's
own second moment.
Using a witness loss as a feasible-action upper bound
requires evaluation under the same loss and membership in
the same exact orthogonal-action class.
Finally, the finite-bank bounds make no population guarantee
for unseen templates.
The separate finite-context certificate requires sampling
assumptions not met by the fixed panel
(\cref{sp:certificate}).
